We evaluated whether sparse feature selection from 4-channel resting-state EEG could identify remission predictors in a home-based tDCS trial for major depression. Across 150 configurations of a hybrid CNN-LSTM and linear SVM, we varied window length (2--10~s) and feature budget (1--70) under leave-one-participant-out cross-validation with SMOTE.

The strongest hybrid configuration (6~s windows, one feature per fold) achieved patient-level ROC-AUC~0.816, PR-AUC~0.579, balanced accuracy~0.786, and MCC~0.571 over 21 held-out participants. A permutation test confirmed above-chance discrimination (raw $p=0.008$; 2-family Bonferroni-corrected $p_\text{corr}=0.016$), although significance does not survive correction for the full 150-configuration sweep. The best SVM (8~s, 30 features) reached ROC-AUC~0.510 ($p=0.469$, uncorrected). The hybrid model exhibited extreme sparsity: only 7 unique features appeared across folds, with modest overlap (mean Jaccard~0.190) but perfect effect-direction consistency among recurrent features.

These findings suggest that patient-level discrimination signals are recoverable from highly sparse representations of low-montage EEG. The recurrent features centred on temporal spectral entropy and frontal-temporal difference measures, converging with prior analyses in the same cohort. The results support low-montage resting-state EEG as a basis for pretreatment stratification research, pending independent validation.
